
\documentclass{article}
\usepackage{colortbl}
\usepackage{makecell}
\usepackage{multirow}
\usepackage{supertabular}

\begin{document}

\newcounter{utterance}

\centering \large Interaction Transcript for game `cladder', experiment `full\_v1.5\_default', episode 9125 with Qwen3.5{-}4B.
\vspace{24pt}

{ \footnotesize  \setcounter{utterance}{1}
\setlength{\tabcolsep}{0pt}
\begin{supertabular}{c@{$\;$}|p{.15\linewidth}@{}p{.15\linewidth}p{.15\linewidth}p{.15\linewidth}p{.15\linewidth}p{.15\linewidth}}
    \# & \multicolumn{2}{c}{Player} && \multicolumn{2}{c}{Game Master} \\
    \hline

    \theutterance \stepcounter{utterance}  
    & & & \multicolumn{4}{p{0.6\linewidth}}{
        \cellcolor[rgb]{0.9,0.9,0.9}{
            \makecell[{{p{\linewidth}}}]{
                \texttt{\tiny{[P1$\langle$GM]}}
                \texttt{You are an expert in causal inference. The following question is not a typical commonsense query, but rather a meticulously designed question created by a professor specializing in causal inference, intended to assess the students' mastery of the course content.} \\
\texttt{Imagine a self{-}contained, hypothetical world with only the following conditions, and without any unmentioned factors or causal relationships: Kraz has a direct effect on pexu and rukz. Pexu has a direct effect on hwax. Hwax has a direct effect on rukz. Kraz is unobserved. For those who are not pexu and are not hwax, the probability of rukz is 19\%. For those who are not pexu and are hwax, the probability of rukz is 29\%. For those who are pexu and are not hwax, the probability of rukz is 20\%. For those who are pexu and are hwax, the probability of rukz is 31\%. For those who are not pexu, the probability of hwax is 16\%. For those who are pexu, the probability of hwax is 90\%. For those who are pexu, would it be less likely to see rukz if the individual was not pexu?} \\
\texttt{Start your answer with 'Yes' or 'No', followed by additional reasoning or evidence to support your explanation.} \\
            }
        }
    }
     \\ \\

    \theutterance \stepcounter{utterance}  
    & \multicolumn{4}{p{0.6\linewidth}}{
        \cellcolor[rgb]{0.561,0.737,0.561}{
            \makecell[{{p{\linewidth}}}]{
                \texttt{\tiny{[P1$\rangle$GM]}}
                \texttt{Yes} \\
\\ 
\texttt{**Reasoning:**} \\
\\ 
\texttt{To answer this question, we must determine the causal effect of the variable **pexu** on the outcome **rukz** by comparing the counterfactual scenarios: the probability of rukz if an individual *is* pexu versus the probability of rukz if that same individual *were not* pexu.} \\
\\ 
\texttt{In this hypothetical world, the causal structure is defined as:} \\
\texttt{1.  \$Kraz \textbackslash{}rightarrow pexu\$} \\
\texttt{2.  \$pexu \textbackslash{}rightarrow hwax\$} \\
\texttt{3.  \$hwax \textbackslash{}rightarrow rukz\$} \\
\texttt{4.  \$Kraz \textbackslash{}rightarrow rukz\$ (direct effect)} \\
\\ 
\texttt{The question asks: "For those who are pexu, would it be less likely to see rukz if the individual was not pexu?"} \\
\texttt{Mathematically, this asks if \$P(rukz | do(pexu=1)) < P(rukz | do(pexu=0))\$.} \\
\\ 
\texttt{We can calculate these probabilities using the provided conditional probabilities and the rules of causal inference (specifically the backdoor adjustment or the causal effect formula).} \\
\\ 
\texttt{**Step 1: Calculate \$P(rukz | do(pexu=0))\$**} \\
\texttt{If we intervene to set \$pexu = 0\$, the probability of \$pexu\$ is fixed at 0. Consequently, the probability of \$hwax\$ is determined solely by the condition "not pexu":} \\
\texttt{*   \$P(hwax | pexu=0) = 16\textbackslash{}\% = 0.16\$} \\
\texttt{*   \$P(rukz | pexu=0, hwax=1) = 29\textbackslash{}\%\$} \\
\texttt{*   \$P(rukz | pexu=0, hwax=0) = 19\textbackslash{}\%\$} \\
\\ 
\texttt{Using the law of total probability:} \\
\texttt{\$\$P(rukz | do(pexu=0)) = P(rukz | pexu=0, hwax=1) \textbackslash{}times P(hwax | pexu=0) + P(rukz | pexu=0, hwax=0) \textbackslash{}times (1 {-} P(hwax | pexu=0))\$\$} \\
\texttt{\$\$P(rukz | do(pexu=0)) = (0.29 \textbackslash{}times 0.16) + (0.19 \textbackslash{}times 0.84)\$\$} \\
\texttt{\$\$P(rukz | do(pexu=0)) = 0.0464 + 0.1596 = 0.206\$\$} \\
\\ 
\texttt{**Step 2: Calculate \$P(rukz | do(pexu=1))\$**} \\
\texttt{If we intervene to set \$pexu = 1\$, the probability of \$pexu\$ is fixed} \\
            }
        }
    }
    & & \\ \\

    \theutterance \stepcounter{utterance}  
    & & & \multicolumn{2}{p{0.3\linewidth}}{
        \cellcolor[rgb]{0.9,0.9,0.9}{
            \makecell[{{p{\linewidth}}}]{
                \texttt{\tiny{[GM$|$GM]}}
                \texttt{Yes} \\
            }
        }
    }
    & & \\ \\

    \theutterance \stepcounter{utterance}  
    & & & \multicolumn{2}{p{0.3\linewidth}}{
        \cellcolor[rgb]{0.9,0.9,0.9}{
            \makecell[{{p{\linewidth}}}]{
                \texttt{\tiny{[GM$|$GM]}}
                \texttt{yes} \\
            }
        }
    }
    & & \\ \\

    \theutterance \stepcounter{utterance}  
    & & & \multicolumn{2}{p{0.3\linewidth}}{
        \cellcolor[rgb]{0.9,0.9,0.9}{
            \makecell[{{p{\linewidth}}}]{
                \texttt{\tiny{[GM$|$GM]}}
                \texttt{game\_result = WIN} \\
            }
        }
    }
    & & \\ \\

\end{supertabular}
}

\end{document}
